% ===== 原文 pp.277–291（PDF p.292–306）· tr-05 · c06c =====
% 本块 = 第 6 章 §6.2 末（从 §6.2.5 的 Glue(𝒮,k) 范畴中间起）+ §6.3 开头（含 6.3.1）；
% tag 6.3–6.5；图 Figure 12 已裁为 work/figures/f6_12.png。不写章标题。

% [原文 p.277]
\begin{itemize}[leftmargin=1.6em,itemsep=2pt]
\item 任意两个对象之间至多有一个态射，并且恰有一个态射
  \[
    (\omega,\tau,u)\to(\omega',\tau',u')
  \]
  当且仅当 $\omega\subseteq\omega'$ 且 $\tau\supseteq\tau'$。
\end{itemize}

注意，该范畴中任一态射都可分解为如下两种基本类型的态射之复合：
\begin{itemize}[leftmargin=1.6em,itemsep=2pt]
\item[(I)] $\omega\subseteq\omega'$，$\tau\supseteq\tau'$，$u=u'$（\emph{换层}）。
\item[(II)] $\omega=\omega'$，$\tau=\tau'$，$\dim u\cap u'=1$，$\omega\cap u\cap u'\ne\varnothing$（\emph{换腔}）。
\end{itemize}

% [Tutor Note: 范畴 $\operatorname{Glue}(\cS,k)$ 的对象是三元组 $(\omega,\tau,u)$，其中 $u$ 恒为粗体（表示腔），而 $\omega,\tau$ 是 $\cP$ 的胞腔；两类基本态射 (I) 换层、(II) 换腔构成了 §6.2.5 全部函子性讨论的骨架。核对锚点：(II) 的条件 $\omega\cap u\cap u'\ne\varnothing$ 与本页第二处出现的形式一致（下一块 p.278 亦同）。]

我们希望定义一个函子
\[
  F_k:\operatorname{Glue}(\cS,k)\to\operatorname{Rings},
\]
从 $\operatorname{Glue}(\cS,k)$ 到环的范畴。对结构 $\cS$ 的任意选取，这个函子都不会是良定义的；但若 $\cS$ 满足一个我们将称之为\emph{相容性}的条件，函子就成为良定义的，我们就可以用它来定义所需的粘合修正。

首先，对 $\operatorname{Glue}(\cS,k)$ 的任一对象 $(\omega,\tau,u)$，如下联系一个环 $R^k_{\omega,\tau,u}$：
\begin{itemize}[leftmargin=1.6em,itemsep=2pt]
\item 若 $\omega=\tau$ 是 $\cP$ 中包含奇点 $p\in\Delta$ 的一条边，则令
  \[
    R^k_{\omega,\tau,u}:=\bigl(R^k_{\omega,\tau,\sigma_u}\bigr)_{f_{\mathfrak d_{p,\pm}}}.
  \]
  这里右端的下标表示在元素 $f_{\mathfrak d_{p,\pm}}=1+z^{m_\pm}$ 处作局部化。注意 $z^{m_\pm}$ 在 $R^k_{\omega,\tau,\sigma_u}$ 中实际上是可逆的，且 $f_{\mathfrak d_{p,+}}=z^{m_-}f_{\mathfrak d_{p,-}}$。因此无论在 $f_{\mathfrak d_{p,+}}$ 处还是 $f_{\mathfrak d_{p,-}}$ 处作局部化，结果都一样。
\item 若 $\omega$ 是顶点、$\tau$ 是 $\cP$ 中包含奇点 $p\in\Delta$ 的边，且取 $\mathfrak d\in\{\mathfrak d_{p,+},\mathfrak d_{p,-}\}$ 为该射线的选取，使得 $\omega$ 是这条射线所遇到的第一条 $\cP$ 的顶点，则令
  \[
    R^k_{\omega,\tau,u}:=\bigl(R^k_{\omega,\tau,\sigma_u}\bigr)_{f_\mathfrak d}.
  \]
\item 其他情形，令
  \[
    R^k_{\omega,\tau,u}:=R^k_{\omega,\tau,\sigma_u}.
  \]
\end{itemize}

我们定义
\[
  F_k(\omega,\tau,u):=R^k_{\omega,\tau,u}.
\]
注意 $R^k_{\omega,\tau,u}$ 是 $R_k=\Bbbk[t]/(t^{k+1})$-代数。

% [Tutor Note: 上面两种局部化的等价性依赖 $f_{\mathfrak d_{p,+}}=z^{m_-}f_{\mathfrak d_{p,-}}$：两个局部化元相差一个可逆元 $z^{m_-}$，故给出同一环。核对锚点：$f_{\mathfrak d_{p,\pm}}=1+z^{m_\pm}$，而 $z^{m_\pm}$ 在 $R^k_{\omega,\tau,\sigma_u}$ 中可逆（见本页第一段的断言）。]

现在希望把 $\operatorname{Glue}(\cS,k)$ 中的每个态射与相应环之间的同态联系起来。我们先对两种基本情形来做：换层与换腔。

I) \emph{换层映射。}设有
\[
  (\omega,\tau,u)\to(\omega',\tau',u)
\]
为 $\operatorname{Glue}(\cS,k)$ 中的一个态射。于是有复合
\[
  R^k_{\omega,\tau,\sigma_u}\twoheadrightarrow R^k_{\omega,\tau',\sigma_u}
  \hookrightarrow R^k_{\omega',\tau',\sigma_u},
\]
它由 $I^k_{\omega,\tau,\sigma_u}\subseteq I^k_{\omega',\tau',\sigma_u}$ 诱导的典范满射与由 $P_{\varphi,\omega,\sigma_u}\subseteq P_{\varphi,\omega',\sigma_u}$ 诱导的典范单射复合而成。在作了适当的局部化之后，这诱导出一个映射，我们下面就来验证。

事实上，考虑第一个映射的相关局部化。这里 $R^k_{\omega,\tau,u}$ 是 $R^k_{\omega,\tau,\sigma_u}$ 在某个元素 $f=1+z^{m_\pm}$ 处的局部化，当且仅当 $\tau$ 是包含奇点的边。

% [原文 p.278]
若 $\tau=\tau'$，则环 $R^k_{\omega,\tau',u}$ 可写成 $R^k_{\omega,\tau',\sigma_u}$ 在同一元素 $f$ 处的局部化。若 $\tau'$ 换成一个顶点，则容易验证 $f$ 在 $R^k_{\omega,\tau',\sigma_u}$ 中可逆。事实上，此时 $\omega$ 必为顶点。若 $\omega\not\subseteq\partial B$，则由引理 6.19 的 (2)，存在 $\sigma\in\cP_{\max}$ 包含 $\tau'=\omega$，使 $\operatorname{ord}_\sigma(m_\pm)>0$。另一方面，若 $\omega\subseteq\partial B$，则存在 $\sigma\in\cP^0_{\max}$ 使 $\operatorname{ord}_\sigma(m_\pm)>0$，因为 $\bar m_\pm$ 不可能与边界相切。于是在任何情形下，$z^{m_\pm}$ 的某个幂落在 $I^k_{\omega,\tau',\sigma_u}$ 中，故 $1+z^{m_\pm}$ 可逆。因此映射 $R^k_{\omega,\tau,u}\to R^k_{\omega,\tau',u}$ 有定义。

其次考虑第二个映射的相关局部化。若 $R^k_{\omega,\tau',u}$ 是 $R^k_{\omega,\tau',\sigma_u}$ 的局部化，则 $R^k_{\omega',\tau',u}$ 是在同一元素处的局部化。于是这些映射在局部化之后有意义，并给出 $R_k$-代数同态
\[
  \psi_{(\omega,\tau),(\omega',\tau')}:R^k_{\omega,\tau,u}\to R^k_{\omega',\tau',u}.
\]
这就是换层映射。

II) \emph{换腔。}设 $u,u'\in\operatorname{Chambers}(\cS,k)$ 满足
\[
  \dim u\cap u'=1\quad\text{且}\quad \omega\cap u\cap u'\ne\varnothing,
\]
并且有 $\operatorname{Glue}(\cS,k)$ 中的态射
\[
  (\omega,\tau,u)\to(\omega,\tau,u').
\]
取一点 $y\in\operatorname{Int}(u\cap u')$，它落在 $(u\cap u')\setminus\Delta$ 的一个与 $\omega$ 相交的连通分支中。特别地，$y$ 不是 $B$ 的奇点。存在唯一的 $n\in\Lambda_y$，它是本原的、零化 $u\cap u'$ 的切空间，且在指向 $u'$ 的切向量上取负值。

对每一对 $(\mathfrak d,x)$，其中 $\mathfrak d\in\cS[k]$、$x\in[0,N^k_\mathfrak d]$ 且 $\mathfrak d(x)=y$，我们得到一个多项式
\[
  f_{(\mathfrak d,x)}=1+c_\mathfrak d z^{m_{\mathfrak d,x}},
\]
其中 $m_{\mathfrak d,x}\in P_{\varphi,\mathfrak d(x)}$ 是 $m_\mathfrak d$ 在 $x$ 处的茎。事实上 $m_{\mathfrak d,x}\in P_{\varphi,\omega,\sigma_u}$。确实，由 $y$ 的选取，存在 $x'\in I_\mathfrak d$ 在 $x$ 附近，使 $\mathfrak d(x')\in\omega$，且 $m_{\mathfrak d,x}$ 与 $m_{\mathfrak d,x'}$ 在平行移动下一致。由于按射线的定义 $m_{\mathfrak d,x'}\in P_{\varphi,\mathfrak d(x')}$，并沿 $\mathfrak d$ 作平行移动得 $P_{\varphi,\mathfrak d(x')}\subseteq P_{\varphi,\omega,\sigma_u}$，故 $m_{\mathfrak d,x}\in P_{\varphi,\omega,\sigma_u}$。于是 $f_{(\mathfrak d,x)}\in R^k_{\omega,\tau,u}$。

现在我们考虑两种情形，分别定义 $R_k$-代数同态 $\theta_{u,u',y}$。

(1) $\sigma_u=\sigma_{u'}$。此时定义同态
\[
  \theta_{u,u',y}:R^k_{\omega,\tau,u}\to R^k_{\omega,\tau,u'}
\]
为
\[
  \theta_{u,u',y}(z^m)=z^m\prod_{(\mathfrak d,x)}f_{(\mathfrak d,x)}^{\langle n,\bar m\rangle}.
\]
注意即使 $\langle n,\bar m\rangle<0$，这也说得通。事实上，由于 $\sigma_u=\sigma_{u'}$，必须有 $u\cap u'\cap\operatorname{Int}(\sigma_u)\ne\varnothing$。因此该乘积中出现的任何 $(\mathfrak d,x)$ 都不满足 $\mathfrak d\in\operatorname{Initial}(\cS)$ 且 $x\in I^{\mathrm{init}}_\mathfrak d$。于是由引理 6.19 的 (2)（当 $\mathfrak d$ 是初始射线时）与定义 6.22 的 (3)（当 $\mathfrak d$ 不是初始射线时），有 $\operatorname{ord}_{\sigma_u}(m_{\mathfrak d,x})>0$。由此推出该乘积中出现的每个 $f_{(\mathfrak d,x)}$ 都可逆。

% [原文 p.279]
\begin{figure}[htbp]
\centering
\includegraphics[width=0.6\textwidth]{f6_12.png}
\figtitle{12}{记号 $\times$ 标出的是奇点 $p$。}
\end{figure}

(2) $\sigma_u\ne\sigma_{u'}$，于是 $\sigma_u\cap\sigma_{u'}=\rho$ 是一条边。此时必有 $\tau\subseteq\rho$。我们定义
\[
  \theta_{u,u',y}:R^k_{\omega,\tau,u}\to R^k_{\omega,\tau,u'}
\]
为如下映射的复合：由平行移动穿过 $y$ 对 $P_{\varphi,\omega,\sigma_u}$ 与 $P_{\varphi,\omega,\sigma_{u'}}$ 作等同所得的映射 $R^k_{\omega,\tau,u}\to R^k_{\omega,\tau,u'}$，与由
\[
  z^m\mapsto z^m\prod_{(\mathfrak d,x)}f_{(\mathfrak d,x)}^{\langle n,\bar m\rangle}
\]
定义的映射 $R^k_{\omega,\tau,u'}\to R^k_{\omega,\tau,u'}$ 的复合。

同样，这说得通，因为 $f_{(\mathfrak d,x)}$ 总是可逆的。事实上，若 $\mathfrak d\in\operatorname{Initial}(\cS)$ 且 $x\in I^{\mathrm{init}}_\mathfrak d$，则 $f_{(\mathfrak d,x)}$ 可逆：或者按环 $R^k_{\omega,\tau,u'}$ 的定义（当 $\tau=\rho$ 时），或者按换层映射讨论中给出的论证（当 $\tau\subsetneq\rho$ 时）。否则，若 $\mathfrak d\notin\operatorname{Initial}(\cS)$，则由于 $y$ 恰好含于两个极大胞腔 $\sigma_\pm$ 中，且 $\bar m_{\mathfrak d,x}$ 平行于 $\rho$，有 $\operatorname{ord}_{\sigma_+}m_{\mathfrak d,x}=\operatorname{ord}_{\sigma_-}m_{\mathfrak d,x}$。于是如前一种情形那样可得 $\operatorname{ord}_{\sigma_\pm}m_{\mathfrak d,x}>0$，故 $f_{(\mathfrak d,x)}$ 在 $R^k_{\omega,\tau,u'}$ 中可逆。

注意，先验地，映射 $\theta_{u,u',y}$ 依赖于点 $y$ 的选取。然而在第一种情形，乘积 $\prod_{(\mathfrak d,x)}f_{(\mathfrak d,x)}\in\Bbbk[P_{\varphi,\omega,\sigma_u}]$ 与 $y$ 的选取无关（由定义 6.24 的 (3)）。因此 $\theta_{u,u',y}$ 的定义与 $y$ 的选取无关，我们可以略去下标 $y$。

在第二种情形，若存在奇点 $p\in\rho$，则同一乘积依赖于 $y$。特别地，$\theta_{u,u',y}$ 似乎会依赖于 $y\in\mathfrak d_{p,+}(I^{\mathrm{init}}_{\mathfrak d_{p,+}})$ 还是 $y\in\mathfrak d_{p,-}(I^{\mathrm{init}}_{\mathfrak d_{p,-}})$。然而事实上 $\theta_{u,u',y}$ 与 $y$ 无关，我们下面就来验证，其计算与 §6.2.3 中所作的计算实质相同。

设 $p\in u\cap u'$，并取 $y_\pm\in u\cap u'\cap\mathfrak d_{p,\pm}(\operatorname{Int}(I^{\mathrm{init}}_{\mathfrak d_{p,\pm}}))$；见图 12。我们利用沿 $y_+$ 从 $\sigma_u$ 到 $\sigma_{u'}$ 的平行移动来等同 $P_{\varphi,\tau,\sigma_u}$ 与 $P_{\varphi,\tau,\sigma_{u'}}$。在这一等同下，通过沿 $y_-$ 从 $\sigma_u$ 到 $\sigma_{u'}$ 的平行移动所得的这两个幺半群的等同，正是 $\cP_\varphi$ 中由绕奇点 $p$ 的逆时针回路给出的单值变换。

% [Tutor Note: 上段“$\theta_{u,u',y}$ 与 $y$ 无关”的要点：沿 $y_+$ 与 $y_-$ 的两条平行移动相差绕奇点 $p$ 的单值变换 $m\mapsto m+\langle\vec d_\omega,\bar m\rangle m_-$（注 6.14），而 $m_+=-m_-$ 使两个乘积逐项相等；这正是 $f_{\mathfrak d_{p,+}}=1+z^{m_+}$ 与 $z^{m_-}+1=z^{m_-}(1+z^{m_+})$ 之差可被吸收的原因。核对锚点：图 12 中 $y_\pm$ 分居 $\rho$ 两侧、$\times$ 标出 $p$。]

% [原文 p.280]
$\cP_\varphi$ 上的这一单值变换由
\[
  m\mapsto m+\langle \vec d_\omega,\bar m\rangle m_-
\]
给出（见注 6.14），其中 $\vec d_\omega$ 是本原的、零化 $\omega$ 的切空间，且在 $\sigma_u$ 上取正值。于是映射 $\theta_{u,u',y_+}$ 由
\[
  \begin{aligned}
    z^m&\mapsto z^mf_{\mathfrak d_{p,+}}^{\langle \vec d_\omega,\bar m\rangle}
    \prod_{(\mathfrak d,x)}f_{(\mathfrak d,x)}^{\langle \vec d_\omega,\bar m\rangle}\\
    &=z^m(1+z^{m_-})^{\langle \vec d_\omega,\bar m\rangle}
    \prod_{(\mathfrak d,x)}f_{(\mathfrak d,x)}^{\langle \vec d_\omega,\bar m\rangle}
  \end{aligned}
\]
给出，其中乘积取遍所有 $(\mathfrak d,x)$，但通常\emph{排除}使 $x\in I^{\mathrm{init}}_{\mathfrak d_{p,+}}$、$\mathfrak d_{p,+}(x)=y_+$ 的那一对 $(\mathfrak d_{p,+},x)$。另一方面，映射 $\theta_{u,u',y_-}$ 由复合
\[
  \begin{aligned}
    z^m&\mapsto z^{m+\langle \vec d_\omega,\bar m\rangle m_-}\\
    &\mapsto z^{m+\langle \vec d_\omega,\bar m\rangle m_-}(1+z^{m_+})^{\langle \vec d_\omega,\bar m\rangle}
    \prod_{(\mathfrak d,x)}f_{(\mathfrak d,x)}^{\langle \vec d_\omega,\bar m\rangle}\\
    &=z^m(z^{m_-}+1)^{\langle \vec d_\omega,\bar m\rangle}
    \prod_{(\mathfrak d,x)}f_{(\mathfrak d,x)}^{\langle \vec d_\omega,\bar m\rangle}
  \end{aligned}
\]
给出，其中记住 $m_+=-m_-$。于是我们看到，此情形下 $\theta_{u,u',y}$ 同样与 $y$ 的选取无关，我们略去 $y$。

我们还没有从 $\operatorname{Glue}(\cS,k)$ 范畴构造出到环范畴的函子。我们只是对 $\operatorname{Glue}(\cS,k)$ 的每个对象联系了一个环，并对每个换层或换腔同态联系了一个环同态。更一般地，给定唯一的态射
\[
  \epsilon:(\omega,\tau,u)\to(\omega',\tau',u'),
\]
我们希望定义
\[
  F_k(\epsilon):R^k_{\omega,\tau,u}\to R^k_{\omega',\tau',u'}.
\]
为此选取一个序列
\[
  u=u_1,u_2,\ldots,u_n=u',
\]
使得
\[
  u_i\cap\omega'\ne\varnothing,\quad \tau'\subseteq\sigma_{u_i},
  \qquad 1\le i\le n,
\]
并且
\[
  \dim u_{i-1}\cap u_i=1,\quad
  u_{i-1}\cap u_i\cap\omega'\ne\varnothing,
  \qquad 2\le i\le n.
\]
于是得到一列换腔映射
\[
  \theta_{u_{i-1},u_i}:R^k_{\omega',\tau',u_{i-1}}\to R^k_{\omega',\tau',u_i},
  \qquad 2\le i\le n,
\]
以及一个换层映射
\[
  \psi_{(\omega,\tau),(\omega',\tau')}:R^k_{\omega,\tau,u_1}\to R^k_{\omega',\tau',u_1}.
\]
我们取
\[
  F_k(\epsilon):=\theta_{u_{n-1},u_n}\circ\cdots\circ\theta_{u_1,u_2}
  \circ\psi_{(\omega,\tau),(\omega',\tau')}.
\]
但仍有问题：这个复合依赖于序列 $\{u_i\}$ 的具体选取。为了使我们将要给出的粘合构造可行，需要

% [原文 p.281]
$F_k(\epsilon)$ 的定义与任何选取无关。为此，我们需要对所考虑的结构再施加一个条件，我们称之为\emph{相容性}。

\begin{definition}
用 $\operatorname{Joints}(\cS,k)$ 记 $\cP_k$ 中不含于 $\partial B$ 的顶点之集，称该集合的元素为\emph{接点}。对 $j$ 是接点，$\sigma_j$ 记 $\cP$ 中含 $j$ 的最小胞腔。
\end{definition}

\begin{definition}
设 $j\in\operatorname{Joints}(\cS,k)$ 是接点。称结构 $\cS$ 在 $j$ 处\emph{到 $k$ 阶是相合的}，如果对含 $j$ 的腔的任一循环排序 $u_1,\ldots,u_n$，复合
\[
  \theta_{u_n,u_1}\circ\theta_{u_{n-1},u_n}\circ\cdots\circ\theta_{u_1,u_2}:
  R^k_{\sigma_j,\sigma_j,u_1}\to R^k_{\sigma_j,\sigma_j,u_1}
\]
是恒等映射。

若 $\cS$ 在每个接点 $j\in\operatorname{Joints}(\cS,k)$ 处都到 $k$ 阶相合，则称 $\cS$ \emph{到 $k$ 阶相容}。
\end{definition}

相容性足以证明态射 $F_k(\epsilon)$ 的良定义性。

\begin{theorem}
设 $\cS$ 到 $k$ 阶相容。

(1) 给定 $\epsilon:(\omega,\tau,u)\to(\omega',\tau',u')$，同态 $F_k(\epsilon)$ 良定义，且与 $u=u_1,\ldots,u_n=u'$ 的选取无关。

(2) 给定 $\epsilon:(\omega,\tau,u)\to(\omega',\tau',u')$ 与
$\epsilon':(\omega',\tau',u')\to(\omega'',\tau'',u'')$，有
\[
  F_k(\epsilon'\circ\epsilon)=F_k(\epsilon')\circ F_k(\epsilon).
\]
于是 $F_k$ 是一个函子。
\end{theorem}

\para{证明。}令
\[
  A:=\{u\in\operatorname{Chambers}(\cS,k)\mid \omega'\cap u\ne\varnothing,\ \tau'\subseteq\sigma_u\}.
\]
定义 $\Sigma$ 为以 $A$ 为顶点集的抽象二维胞腔复形，当 $\dim u\cap u'=1$ 时把 $u,u'\in A$ 用边相连，并对任一具有公共接点的腔的循环 $u_1,\ldots,u_n$ 粘上一个圆盘。注意若存在这样的接点，则 $\omega'=\tau'=\sigma_j$。一条以 $u,u'$ 为顶点的边给出换腔映射 $\theta_{u,u'}:R^k_{\omega',\tau',u}\to R^k_{\omega',\tau',u'}$。相合性说明沿二维胞腔边界作这些映射的复合为恒等。于是只要 $\Sigma$ 单连通，就得到 (1) 中所需的无关性。

为看出 $\Sigma$ 单连通，我们如下进行。令
\[
  V:=\bigcup_{\substack{\sigma\in\cP\\ \tau'\subseteq\sigma}}\operatorname{Int}(\sigma)
\]
记 $\tau'$ 关于 $\cP$ 的开星。则对 $u\in\operatorname{Chambers}(\cS,k)$，条件 $\tau'\subseteq\sigma_u$ 等价于 $\operatorname{Int}(u)\subseteq V$，而这样的腔把 $V$ 分解为多面体。称该分解为 $\cP_V$。由于 $\omega'=\tau'$ 在拓扑上是一个球，$\cP_V$ 中与 $\omega'$ 相交的胞腔构成一个二维球的多面体分解 $\cP'$。现在 $\Sigma$ 是 $\cP'$ 的对偶多面体分解，从而是单连通的。

至于 (2)，设 $u_1,\ldots,u_n$ 是用于计算 $F_k(\epsilon)$ 的一列腔，而 $u_n=u'_1,\ldots,u'_m$ 是用于计算 $F_k(\epsilon')$ 的一列腔，于是
\[
  \begin{aligned}
    F_k(\epsilon')\circ F_k(\epsilon)
    &=\theta_{u'_{m-1},u'_m}\circ\cdots\circ\theta_{u'_1,u'_2}
    \circ\psi_{(\omega',\tau'),(\omega'',\tau'')}\\
    &\quad\circ\theta_{u_{n-1},u_n}\circ\cdots\circ\theta_{u_1,u_2}
    \circ\psi_{(\omega,\tau),(\omega',\tau')}.
  \end{aligned}
\]

% [Tutor Note: 证明中 $\Sigma$ 的单连通性通过“对偶多面体分解”得到：$\cP_V$ 中与 $\omega'$ 相交的胞腔给出二维球的多面体分解 $\cP'$，其单纯对偶 $\Sigma$ 单连通。核对锚点：$A$ 的顶点是满足 $\omega'\cap u\ne\varnothing$、$\tau'\subseteq\sigma_u$ 的腔；这也是 $F_k(\epsilon)$ 的定义能沿任意腔序列一致的原因。]

% [原文 p.282]
容易验证 $\psi_{(\omega',\tau'),(\omega'',\tau'')}$ 与每个 $\theta_{u_{i-1},u_i}$（$2\le i\le n$）都可交换，只要注意 $\theta_{u_{i-1},u_i}$ 也可解释为映射
\[
  R^k_{\omega'',\tau'',u_{i-1}}\to R^k_{\omega'',\tau'',u_i},
\]
因为
\[
  u_{i-1}\cap u_i\cap\omega'\ne\varnothing
  \Longrightarrow u_{i-1}\cap u_i\cap\omega''\ne\varnothing,
\]
\[
  \tau'\subseteq\sigma_{u_i}\Longrightarrow\tau''\subseteq\sigma_{u_i}.
\]
于是
\[
  \begin{aligned}
    F_k(\epsilon')\circ F_k(\epsilon)
    &=\theta_{u'_{m-1},u'_m}\circ\cdots\circ\theta_{u_1,u_2}
    \circ\psi_{(\omega',\tau'),(\omega'',\tau'')}\circ\psi_{(\omega,\tau),(\omega',\tau')}\\
    &=\theta_{u'_{m-1},u'_m}\circ\cdots\circ\theta_{u_1,u_2}
    \circ\psi_{(\omega,\tau),(\omega'',\tau'')}\\
    &=F_k(\epsilon'\circ\epsilon),
  \end{aligned}
\]
如所断言，其中用序列 $u_1,\ldots,u_n=u'_1,\ldots,u'_m$ 来定义 $F_k(\epsilon'\circ\epsilon)$。

综上，给定到 $k$ 阶相容的结构 $\cS$，我们构造出函子
\[
  F_k:\operatorname{Glue}(\cS,k)\to\operatorname{Rings},
\]
其中
\[
  F_k(\omega,\tau,u)=R^k_{\omega,\tau,u}.
\]
事实上，$F_k$ 取值于 $R_k$-代数的范畴。

\begin{example}
回到例 6.23，那里给出的结构 $\cS$ 对一切阶都相容。唯一的接点在原点，取 $u_i=\sigma_i$，可算得
\[
  \begin{aligned}
    &\theta_{u_5,u_1}\circ\cdots\circ\theta_{u_1,u_2}(z^m)\\
    &=\theta_{u_5,u_1}\circ\cdots\circ\theta_{u_2,u_3}
    \Bigl(z^m\bigl(1+z^{(0,-1,-1)}\bigr)^{\langle(1,-1),\bar m\rangle}\Bigr)\\
    &=\theta_{u_5,u_1}\circ\cdots\circ\theta_{u_3,u_4}
    \Biggl(z^m\bigl(1+z^{(0,0,-1)}\bigr)^{\langle(1,0),\bar m\rangle}\\
    &\qquad\cdot\Bigl(1+\frac{z^{(0,-1,-1)}}{1+z^{(0,0,-1)}}\Bigr)^{\langle(1,-1),\bar m\rangle}\Biggr)\\
    &=\theta_{u_5,u_1}\circ\theta_{u_4,u_5}
    \Biggl(z^m\bigl(1+z^{(0,-1,0)}\bigr)^{\langle(0,1),\bar m\rangle}
    \Bigl(1+\frac{z^{(0,0,-1)}}{1+z^{(0,-1,0)}}\Bigr)^{\langle(1,0),\bar m\rangle}\\
    &\qquad\cdot\Bigl(1+\frac{z^{(0,-1,-1)}}{((1+z^{(0,0,-1)})/(1+z^{(0,-1,0)}))(1+z^{(0,-1,0)})}\Bigr)^{\langle(1,-1),\bar m\rangle}\Biggr)
  \end{aligned}
\]

% [原文 p.283]
\[
  \begin{aligned}
  &=\theta_{u_5,u_1}\circ\theta_{u_4,u_5}
  \Biggl(z^m\bigl(1+z^{(0,-1,0)}\bigr)^{\langle(0,1),\bar m\rangle}
  \Bigl(1+\frac{z^{(0,0,-1)}}{1+z^{(0,-1,0)}}\Bigr)^{\langle(1,0),\bar m\rangle}\\
  &\qquad\cdot\Bigl(1+\frac{z^{(0,-1,-1)}}{1+z^{(0,-1,0)}+z^{(0,0,-1)}}\Bigr)^{\langle(1,-1),\bar m\rangle}\Biggr)\\
  &=\theta_{u_5,u_1}
  \Biggl(z^m\bigl(1+z^{(0,0,-1)}\bigr)^{\langle(-1,0),\bar m\rangle}
  \bigl(1+z^{(0,-1,0)}(1+z^{(0,0,-1)})\bigr)^{\langle(0,1),\bar m\rangle}\\
  &\qquad\cdot\Bigl(1+\frac{z^{(0,0,-1)}}{1+z^{(0,-1,0)}(1+z^{(0,0,-1)})}\Bigr)^{\langle(1,0),\bar m\rangle}\\
  &\qquad\cdot\Bigl(1+\frac{z^{(0,-1,-1)}(1+z^{(0,0,-1)})}
  {1+z^{(0,-1,0)}(1+z^{(0,0,-1)})+z^{(0,0,-1)}}\Bigr)^{\langle(1,-1),\bar m\rangle}\Biggr)\\
  &=\theta_{u_5,u_1}
  \Biggl(z^m\bigl(1+z^{(0,0,-1)}\bigr)^{\langle(-1,0),\bar m\rangle}
  \bigl(1+z^{(0,-1,0)}+z^{(0,-1,-1)}\bigr)^{\langle(0,1),\bar m\rangle}\\
  &\qquad\cdot\Bigl(\frac{(1+z^{(0,-1,0)})(1+z^{(0,0,-1)})}
  {1+z^{(0,-1,0)}+z^{(0,-1,-1)}}\Bigr)^{\langle(1,0),\bar m\rangle}
  \Bigl(1+\frac{z^{(0,-1,-1)}}{1+z^{(0,-1,0)}}\Bigr)^{\langle(1,-1),\bar m\rangle}\Biggr)\\
  &=\theta_{u_5,u_1}
  \Biggl(z^m\bigl(1+z^{(0,-1,0)}+z^{(0,-1,-1)}\bigr)^{\langle(0,1),\bar m\rangle}
  \Bigl(\frac{1+z^{(0,-1,0)}}{1+z^{(0,-1,0)}+z^{(0,-1,-1)}}\Bigr)^{\langle(1,0),\bar m\rangle}\\
  &\qquad\cdot\Bigl(\frac{1+z^{(0,-1,0)}+z^{(0,-1,-1)}}{1+z^{(0,-1,0)}}\Bigr)^{\langle(1,-1),\bar m\rangle}\Biggr)\\
  &=\theta_{u_5,u_1}\Bigl(z^m\bigl(1+z^{(0,-1,0)}\bigr)^{\langle(0,1),\bar m\rangle}\Bigr)
  =z^m.
  \end{aligned}
\]
这里，散射图中初始射线之外的那条额外对角射线对于保证相合性是至关重要的，因为与穿过该接点的两条初始射线相联系的自治同态并不交换。
\end{example}

\para{6.2.6.\quad 由相容结构得到 $k$ 阶形变。}固定 $B$ 上到 $k$ 阶相容的结构 $\cS$。这给出函子
\[
  F_k:\operatorname{Glue}(\cS,k)\to\operatorname{Rings}.
\]
现在要定义 $\check X_0(B,\cP)$ 的 $k$ 阶形变。做法与 §6.2.2 基本相同，只是把映射换成由 $F_k$ 给出的映射。

首先固定 $\omega\in\cP$。我们将定义与 §6.2.2 中类似的一个概形 $U^k_\omega$。对每个含 $\omega$ 的 $\tau\in\cP$，选取 $u_\tau\in\operatorname{Chambers}(\cS,k)$，使得 $u_\tau\cap\omega\ne\varnothing$、$\tau\subseteq\sigma_{u_\tau}$。于是给定 $\omega\subseteq\tau_1\subseteq\tau_2$，我们得到 $\operatorname{Glue}(\cS,k)$ 中的态射
\[
  \epsilon_{\tau_1,\tau_2}:(\omega,\tau_2,u_{\tau_2})\to(\omega,\tau_1,u_{\tau_1}),
\]
从而得到同态
\[
  F_k(\epsilon_{\tau_1,\tau_2}):R^k_{\omega,\tau_2,u_{\tau_2}}\to R^k_{\omega,\tau_1,u_{\tau_1}}.
\]

% [原文 p.284]
这给出环的逆向系统：由于 $F_k$ 是函子，有
$F_k(\epsilon_{\tau_2,\tau_3})\circ F_k(\epsilon_{\tau_1,\tau_2})=F_k(\epsilon_{\tau_1,\tau_3})$。我们令
\[
  R^k_\omega:=\varprojlim_{\tau\supseteq\omega}R^k_{\omega,\tau,u_\tau}
\]
以及
\[
  U^k_\omega:=\operatorname{Spec}R^k_\omega.
\]
这些定义事实上与诸 $u_\tau$ 的选取无关。确实，给定另一组选取 $u'_\tau$（$\tau\supseteq\omega$），给出映射
$\epsilon'_{\tau_1,\tau_2}:(\omega,\tau,u'_{\tau_2})\to(\omega,\tau,u'_{\tau_1})$，则存在唯一的映射
\[
  \epsilon_\tau:(\omega,\tau,u_\tau)\to(\omega,\tau,u'_\tau),
\]
满足
\[
  F_k(\epsilon'_{\tau_1,\tau_2})\circ F_k(\epsilon_{\tau_2})
  =F_k(\epsilon_{\tau_1})\circ F_k(\epsilon_{\tau_1,\tau_2}).
\]
此外，由于 $F_k(\epsilon_\tau)$ 由一列换腔映射给出，$F_k(\epsilon_\tau)$ 是同构，从而给出诸 $u_\tau$ 与诸 $u'_\tau$ 所定义的逆向系统之间的同构。

\begin{lemma}
$U^0_\omega$ 同构于 $\operatorname{Spec}\Bbbk[P_{\varphi,x}]/(t)$，其中 $x\in\operatorname{Int}(\omega)\cap B_0$。
\end{lemma}

\para{证明。}注意可以取 $\cP^0=\cP$，因为 $\cS[0]$ 可取为仅由 $\cS$ 的初始射线组成，而对初始射线 $\mathfrak d$，可以取 $N^0_\mathfrak d$ 使得 $[0,N^0_\mathfrak d]=I^{\mathrm{init}}_\mathfrak d$。于是可以取 $\operatorname{Chambers}(\cS,0)=\cP_{\max}$。现在固定 $\omega$ 与 $x\in\operatorname{Int}(\omega)\cap B_0$。注意对每个 $(\omega,\tau,u)\in\operatorname{Glue}(\cS,0)$，$R^0_{\omega,\tau,\sigma_u}$ 可以通过从 $\sigma_u$ 的内部到 $x$ 的平行移动，等同于 $\Bbbk[P_{\varphi,x}]/I^0_{\omega,\tau,x}$，其中 $I^0_{\omega,\tau,x}$ 是 $P_{\varphi,x}$ 中位于扇 $\Sigma_x$ 中切锥 $T_x$ 之上的那个面的补集。特别地，$R^0_{\omega,\tau,\sigma_u}$ 是整环，且 $R^0_{\omega,\tau,u}$ 是 $R^0_{\omega,\tau,\sigma_u}$ 的局部化，故 $R^0_{\omega,\tau,\sigma_u}\subseteq R^0_{\omega,\tau,u}$（除非 $\dim\tau=1$ 时取等号）。换层映射
\[
  \epsilon:(\omega,\tau_2,u)\to(\omega,\tau_1,u)
\]
给出
\[
  F_0(\epsilon):R^0_{\omega,\tau_2,u}\to R^0_{\omega,\tau_1,u},
\]
它正是典范满射
\[
  \Bbbk[P_{\varphi,x}]/I^0_{\omega,\tau_2,x}
  \twoheadrightarrow\Bbbk[P_{\varphi,x}]/I^0_{\omega,\tau_1,x}
\]
的局部化。另一方面，唯一可能非平凡的换腔映射形如 $\epsilon:(\omega,\tau,u)\to(\omega,\tau,u')$，其中 $\dim\tau=1$，而 $u,u'$ 是含 $\tau$ 的两个不同的极大胞腔。但若 $z^m\in\Bbbk[P_{\varphi,x}]/I^0_{\omega,\tau,x}$ 非零，则 $\bar m$ 必与 $\tau$ 相切，从而按 $\theta_{u,u'}$ 的定义有 $\theta_{u,u'}(z^m)=z^m$。因此该换腔映射是平凡的。

由此推出逆向系统之间存在包含
\[
  \bigl(\Bbbk[P_{\varphi,x}]/I^0_{\omega,\tau,x}\bigr)_\tau
  \hookrightarrow\bigl(R^0_{\omega,\tau,u_\tau}\bigr)_\tau,
\]
其中左端逆向系统中的同态是典范的，右端的同态正是定义 $R^0_\omega$ 的那些同态。此外，除非 $\dim\tau=1$，$\Bbbk[P_{\varphi,x}]/I^0_{\omega,\tau,x}\cong R^0_{\omega,\tau,u_\tau}$。由于逆向极限 $\varprojlim\Bbbk[P_{\varphi,x}]/I^0_{\omega,\tau,x}$ 中的元素 $(f_\tau)$ 完全由那些 $\tau\in\cP_{\max}$ 对应的 $f_\tau$ 决定，$\varprojlim R^0_{\omega,\tau,u_\tau}$ 也如此，故这两个逆向极限事实上同构。于是
\[
  R^0_\omega\cong\varprojlim\Bbbk[P_{\varphi,x}]/I^0_{\omega,\tau,x}
  \cong\Bbbk[P_{\varphi,x}]/(t),
\]
最后一个同构如引理 6.9 中那样。

% [原文 p.285]
下面将证明 $U^k_\omega$ 是 $U^0_\omega$ 在 $O_k=\operatorname{Spec}\Bbbk[t]/(t^{k+1})$ 上的平坦形变。在开始之前，需要一个观察。由于所有环 $R^k_{\omega,\tau,u}$ 都是 $R_k=\Bbbk[t]/(t^{k+1})$-代数，$R^k_\omega$ 也是 $R_k$-代数。于是对每个 $\ell\le k$，有 $R_k$-模的复形
\begin{equation}
  0\to R^0_\omega\xrightarrow{\ \cdot t^\ell\ }R^\ell_\omega
  \to R^{\ell-1}_\omega\to 0. \tag{6.3}
\end{equation}
这里映射 $R^0_\omega\to R^\ell_\omega$ 由
\[
  \varprojlim R^0_{\omega,\tau,u_\tau}\ni(f_\tau)_\tau
  \mapsto(t^\ell f_\tau)_\tau
\]
给出。由于 $t^\ell$ 零化任一满足 $z^m\in I^0_{\omega,\tau,\sigma_{u_\tau}}$ 的元素 $z^m\in R^\ell_{\omega,\tau,u_\tau}$，可见 $(t^\ell f_\tau)_\tau\in\varprojlim R^\ell_{\omega,\tau,u_\tau}$。映射 $R^\ell_\omega\to R^{\ell-1}_\omega$ 由 $(f_\tau)_\tau\mapsto(f_\tau\bmod I^{\ell-1}_{\omega,\tau,\sigma_{u_\tau}})$ 给出。先验地，(6.3) 只是复形而非正合列。

% [Tutor Note: (6.3) 的正合性是后面平坦性判据的全部内容：由 (6.3) 正合可归纳得到 $R^k_\omega/tR^k_\omega\cong R^0_\omega$，再引用 [78] 定理 22.3 的条件 (4) 得 $R^k_\omega$ 是平坦 $R_k$-代数。核对锚点：$R^\ell_\omega\to R^{\ell-1}_\omega$ 由模 $I^{\ell-1}_{\omega,\tau,\sigma_{u_\tau}}$ 给出。]

\begin{lemma}
设 (6.3) 对每个 $\ell\le k$ 都正合。则 $U^k_\omega$ 是 $U^0_\omega$ 在 $O_k$ 上的平坦形变。
\end{lemma}

\para{证明。}若 (6.3) 正合，则归纳可得由 $(f_\tau)_\tau\mapsto(f_\tau\bmod I^0_{\omega,\tau,\sigma_{u_\tau}})_\tau$ 给出的映射 $R^\ell_\omega\to R^0_\omega$ 是满射。于是特别地，由乘以 $t^\ell$ 给出的映射 $R^0_\omega\to R^\ell_\omega$ 映入 $t^\ell R^\ell_\omega$。然而无论如何，显然该映射的像包含 $t^\ell R^\ell_\omega$。因此 $R^0_\omega$ 在 $R^\ell_\omega$ 中的像（由乘以 $t^\ell$ 给出）就是 $t^\ell R^\ell_\omega$，故 (6.3) 的正合性蕴含 $R^{\ell-1}_\omega\cong R^\ell_\omega/t^\ell R^\ell_\omega$。归纳可得 $R^k_\omega/tR^k_\omega\cong R^0_\omega$。再者，由 [78] 的定理 22.3（应用该定理的条件 (4)），$R^k_\omega$ 是平坦的 $R_k$-代数。于是 $U^k_\omega$ 是 $U^0_\omega$ 在 $O_k$ 上的平坦形变。

下文将用到 $U^0_\omega$（$\omega\in\cP$）的环面层。给定 $\tau\in\cP$ 且 $\tau\supseteq\omega$，有 $U^0_\omega$ 中对应于 $\tau$ 的环面层 $V_\tau$。事实上 $V_\tau\cong\operatorname{Spec}R^0_{\omega,\tau,\sigma_{u_\tau}}$。现在如下定义 $U^0_\omega$ 的一维层之子集 $Z_\omega\subseteq U^0_\omega$。若 $\tau\supseteq\omega$ 是一维胞腔，则当 $\tau\cap\Delta=\varnothing$ 时取 $Z_\omega\cap V_\tau$ 为空集。若 $\tau\cap\Delta=\{p\}$，则 $f_{\mathfrak d_{p,\pm}}$ 之一可看作 $R^0_{\omega,\tau,\sigma_{u_\tau}}$ 的元素，从而生成一个理想。该理想就是 $Z_\omega\cap V_\tau$ 的理想。当然，根据 $f_{\mathfrak d_{p,\pm}}$ 的显式形式，$Z_\omega\cap V_\tau$ 只由一个点组成。然后取 $Z_\omega$ 为这些集合 $Z_\omega\cap V_\tau$ 在全部含 $\omega$ 的一维 $\tau$ 上的并。

现在可以陈述本小节的主要结果。

\begin{theorem}
$U^k_\omega$ 是 $U^0_\omega$ 在 $O_k$ 上的平坦形变。因此 $U^k_\omega$ 与 $U^0_\omega$ 的底层拓扑空间相同。特别地，设 $y\in U^k_\omega$ 作为 $U^0_\omega$ 的点含于某个含 $\omega$ 但不含于任何 $V_{\tau'}\subseteq V_\tau$（$\tau'\subsetneq\tau$）的层 $V_\tau$ 中。则若 $y\notin Z_\omega$，$U^k_\omega$ 在 $y$ 的某邻域内 étale 局部同构于 $\operatorname{Spec}\Bbbk[P_{\varphi,x}]/(t^{k+1})$，其中 $x\in\operatorname{Int}(\tau)\cap B_0$。
\end{theorem}

\para{证明。}有 $0\le\dim\omega\le2$，故对 $\dim\omega$ 有三种情形，我们将逐一处理，从最易到最难。

\para{$\dim\omega=2$：}此时 $\omega\in\cP_{\max}$，$\sigma_{u_\omega}=\omega$，故对 $x\in\operatorname{Int}(\omega)$，$P_{\varphi,\omega,\sigma_{u_\omega}}$ 同构于 $\Lambda_x\times\mathbb{N}$。于是
\[
  R^k_{\omega,\omega,u_\omega}\cong\Bbbk[\Lambda_x]\otimes_\Bbbk R_k,
\]
从而
\[
  U^k_\omega\cong U^0_\omega\times_\Bbbk O_k.
\]
这显然是 $U^0_\omega$ 的平坦形变，且同构于 $\operatorname{Spec}\Bbbk[P_{\varphi,x}]/(t^{k+1})$。

% [原文 p.286]
\para{$\dim\omega=1$。}此时有 $\omega=\sigma_+\cap\sigma_-$，其中 $\sigma_\pm$ 是两个极大胞腔，而定义 $R^k_\omega$ 的逆向极限事实上是纤维积
\[
  R^k_\omega=R^k_{\omega,\sigma_-,u_{\sigma_-}}
  \times_{R^k_{\omega,\omega,u_\omega}}R^k_{\omega,\sigma_+,u_{\sigma_+}}.
\]
我们将在 $u_{\sigma_+}$ 与 $u_{\sigma_-}$ 取为相邻（即 $\dim u_{\sigma_+}\cap u_{\sigma_-}=1$，且 $u_{\sigma_+}$ 与 $u_{\sigma_-}$ 自然位于 $\omega$ 的两侧）的情形下显式计算。也可以取 $u_\omega=u_{\sigma_+}$。此时可以取 $\varphi$ 的一个代表，它在 $\sigma_-$ 上斜率为 $0$，在 $\sigma_+$ 上斜率为 $e\vec d_\omega$，其中 $e$ 是正整数。如通常那样，$\vec d_\omega\in\Lambda_y$（$y\in u_{\sigma_+}\cap u_{\sigma_-}$ 用于定义换腔映射）。该元素是本原的、在 $\omega$ 的切空间上为零，且在指向 $\sigma_+$ 的切向量上取正值。我们可以通过穿过 $y$ 的平行移动把 $P_{\varphi,\omega,\sigma_+}$ 与 $P_{\varphi,\omega,\sigma_-}$ 等同。容易看出
\[
  P_{\varphi,\omega,\sigma_+}\cong\Lambda_\omega\oplus S_e,
\]
其中 $\Lambda_\omega\subseteq\Lambda_y$ 是 $\omega$ 在 $y$ 处的整切向量空间，$S_e$ 是例 3.28 中定义的幺半群。我们把 $S_e\subseteq\Z^2$ 看作由 $(-1,0)$、$(0,1)$、$(1,e)$ 生成的幺半群，从而得到 $\Bbbk[S_e]$ 的元素
\[
  u=z^{(-1,0)},\qquad v=z^{(1,e)},\qquad t=z^{(0,1)}.
\]
满足
\[
  \operatorname{ord}_{\sigma_-}u=0,\quad
  \operatorname{ord}_{\sigma_-}v=e,\quad
  \operatorname{ord}_{\sigma_-}t=1,
\]
\[
  \operatorname{ord}_{\sigma_+}u=e,\quad
  \operatorname{ord}_{\sigma_+}v=0,\quad
  \operatorname{ord}_{\sigma_+}t=1.
\]
记
\[
  R_\pm:=R^k_{\omega,\sigma_\pm,u_{\sigma_\pm}},\qquad
  R_\cap:=R^k_{\omega,\omega,u_\omega},
\]
则
\[
  \begin{aligned}
    R_-&=\Bbbk[\Lambda_\omega][u,v,t]/\langle uv-t^e,\ v^\beta t^\gamma\mid
    \beta e+\gamma\ge k+1\rangle,\\
    R_+&=\Bbbk[\Lambda_\omega][u,v,t]/\langle uv-t^e,\ u^\alpha t^\gamma\mid
    \alpha e+\gamma\ge k+1\rangle,\\
    R_\cap&=\bigl(\Bbbk[\Lambda_\omega][u,v,t]/\langle uv-t^e,\ u^\alpha v^\beta t^\gamma\mid
    \max\{\alpha,\beta\}e+\gamma\ge k+1\rangle\bigr)_{f_{\mathfrak d_{p,\pm}}}.
  \end{aligned}
\]
由于我们取 $u_\omega=u_{\sigma_+}$，映射 $R_+\to R_\cap$ 是典范满射再作局部化。另一方面，映射 $R_-\to R_\cap$ 是典范满射 $R^k_{\omega,\sigma_-,\sigma_-}\to R^k_{\omega,\omega,\sigma_-}$ 再作局部化，然后作换腔映射，而该换腔映射在这些坐标下由
\[
  u\mapsto uf_\omega,\qquad v\mapsto vf_\omega^{-1},\qquad t\mapsto t
\]
给出，其中
\[
  f_\omega=\prod_{(\mathfrak d,x)}f_{(\mathfrak d,x)},
\]
乘积取遍所有满足 $\mathfrak d(x)=y$ 的 $(\mathfrak d,x)$。可以把 $f_\omega$ 看作环 $\Bbbk[\Lambda_\omega][t]$ 的元素。

\begin{lemma}
令
\[
  R_\cup:=\Bbbk[\Lambda_\omega][U,V,t]/(UV-f_\omega t^e,t^{k+1}).
\]

% [原文 p.287]
则
\[
  \begin{aligned}
    R_\cup&\to R_-\times_{R_\cap}R_+\\
    U&\mapsto(u,f_\omega u)\\
    V&\mapsto(f_\omega v,v)\\
    t&\mapsto(t,t)
  \end{aligned}
\]
是 $\Bbbk[\Lambda_\omega][t]$-代数的同构。
\end{lemma}

\para{证明。}注意环 $R_-$ 与 $R_+$ 作为 $\Bbbk[\Lambda_\omega][t]/(t^{k+1})$-模由 $1$、$u^i$、$v^j$（$i,j>0$）生成。又，由 $u^i$（$i\ge0$）与 $v^j$（$j\ge0$）生成的 $R_-$（$R_+$）的子模是自由直和项。于是对 $g_\pm\in R_\pm$，可以唯一地写成
\[
  g_-=\sum_{i\ge0}a_iu^i+h_-(v,t),\qquad
  g_+=\sum_{j\ge0}b_jv^j+h_+(u,t),
\]
其中 $a_i,b_j\in\Bbbk[\Lambda_\omega][t]/(t^{k+1})$、$h_\pm(0,t)=0$。于是 $(g_-,g_+)\in R_-\times_{R_\cap}R_+$ 当且仅当
\[
  a_0=b_0,\qquad
  h_-(v,t)=\sum_{j>0}b_jf_\omega^jv^j,\qquad
  h_+(u,t)=\sum_{i>0}a_if_\omega^iu^i
\]
作为 $R_\cap$ 的元素成立。但若如此，则 $(g_-,g_+)$ 是 $\sum_{i\ge0}a_iU^i+\sum_{j>0}b_jV^j\in R_\cup$ 的像。这说明满射性。单射性容易由注意到 $R_\cup$ 是以 $U^i,V^j$（$i\ge0$，$j>0$）为基的自由 $\Bbbk[\Lambda_\omega][t]/(t^{k+1})$-模而得到。

作为该引理的推论，我们看到 $R^k_\omega=R_\cup$，由此容易推出 (6.3) 正合。这说明 $U^k_\omega$ 是 $U^0_\omega$ 的平坦形变。此外，容易验证，在 $Z_\omega$（即 $\operatorname{Spec}R_\cup$ 中的轨迹 $x=y=f_\omega=t=0$）之外，$U^k_\omega$ 具有所需的局部形式。

% [Tutor Note: $\dim\omega=1$ 情形中，纤维积 $R^k_\omega=R^k_{\omega,\sigma_-,u_{\sigma_-}}\times_{R^k_{\omega,\omega,u_\omega}}R^k_{\omega,\sigma_+,u_{\sigma_+}}$ 的显式计算（引理 6.33）关键在于 $uv=t^e$ 把 $u$ 与 $v$ 的幂次纠缠起来；换腔映射 $u\mapsto uf_\omega,\ v\mapsto vf_\omega^{-1}$ 中的 $f_\omega$ 正好被 $\sum a_iU^i+\sum b_jV^j$ 吸收，故 $R_\cup\to R_-\times_{R_\cap}R_+$ 可逆。核对锚点：$R_\cup$ 的自由基为 $U^i,V^j$（$i\ge0,j>0$）。]

\para{$\dim\omega=0$。}这是最困难的情形，我们分两步处理。固定 $u_\omega$，从而固定幺半群
\[
  P:=P_{\varphi,\omega,\sigma_{u_\omega}}.
\]
则由引理 6.30，
\[
  U^0_\omega=\operatorname{Spec}\Bbbk[P]/(t).
\]
局部地工作是有益的。为此我们先定义一个颇为奇特的对象。首先注意，与 $\operatorname{Spec}R^k_{\omega,\tau,u}$ 相伴的既约概形是 $\operatorname{Spec}R^0_{\omega,\tau,u}$，它是 $\operatorname{Spec}R^0_{\omega,\tau,\sigma_u}$ 的开子概形，而后者是 $U^0_\omega$ 的闭层。对概形 $X$，用 $|X|$ 记其底层拓扑空间。于是有包含
\[
  i_{\omega,\tau}:|\operatorname{Spec}R^k_{\omega,\tau,u_\tau}|\to|U^0_\omega|.
\]
令
\[
  \cO^k_{\omega,\tau}:=i_{\omega,\tau*}\cO_{\operatorname{Spec}R^k_{\omega,\tau,u_\tau}}.
\]
这是 $|U^0_\omega|$ 上的环层。注意映射 $F_k(\epsilon_{\tau_1,\tau_2}):R^k_{\omega,\tau_2,u_{\tau_2}}\to R^k_{\omega,\tau_1,u_{\tau_1}}$ 诱导出映射 $\cO^k_{\omega,\tau_2}\to\cO^k_{\omega,\tau_1}$，从而定义层的逆向系统 $(\cO^k_{\omega,\tau})_\tau$。令
\[
  \cO^k_\omega:=\varprojlim\cO^k_{\omega,\tau},
\]

% [原文 p.288]
并记环化空间
\[
  V^k_\omega:=(|U^0_\omega|,\cO^k_\omega).
\]
与复形 (6.3) 完全类似，对 $\ell\le k$ 有 $|U^0_\omega|$ 上层的复形
\begin{equation}
  0\to\cO^0_\omega\xrightarrow{\ \cdot t^\ell\ }\cO^\ell_\omega
  \to\cO^{\ell-1}_\omega\to 0. \tag{6.4}
\end{equation}
由于层的逆向极限 $\varprojlim\cF_i$ 由 $U\mapsto\varprojlim\cF_i(U)$ 定义，对 (6.4) 取整体截面确实得到复形 (6.3)：这里用到 $\cO^0_\omega\cong\cO_{U^0_\omega}$，它由引理 6.30 的论证可得。由于 $U^0_\omega$ 是仿射概形，有 $H^1(|U^0_\omega|,\cO_{U^0_\omega})=0$，故若 (6.4) 正合，则 (6.3) 正合。

现在考察开集 $|U^0_\tau|\subseteq|U^0_\omega|$，其中 $\tau\in\cP$ 含 $\omega$ 且 $\dim\tau=1$。利用逆向极限的性质容易看出，环化空间 $(|U^0_\tau|,\cO^k_\omega|_{|U^0_\tau|})$ 同构于仿射概形 $U^k_\tau$。我们已经对 $U^k_\tau$ 证明了定理 6.32，因此 (6.4) 在 $\bigcup_{\tau\supseteq\omega,\dim\tau=1}|U^0_\tau|$ 上正合。该集合几乎就是整个 $|U^0_\omega|$：它只漏掉一个点 $z\in|U^0_\omega|$，对应于 $\Bbbk[P]/(t)$ 中由幺半群理想 $P\setminus\{0\}\subseteq P$ 给出的极大理想 $\mathfrak m_0$。用 $\mathfrak m_k$ 记 $\Bbbk[P]/(t^{k+1})$ 中由同一幺半群理想 $P\setminus\{0\}$ 给出的极大理想。

\begin{lemma}
存在同构
\[
  \psi:\bigl(\Bbbk[P]/(t^{k+1})\bigr)_{\mathfrak m_k}\to\cO^k_{\omega,z}.
\]
\end{lemma}

\para{证明。}取有限逆向极限与取茎可交换，故
\[
  \cO^k_{\omega,z}=\varprojlim(\cO^k_{\omega,\tau})_z.
\]
由于 $\cO^k_{\omega,\tau}$ 支在 $z$ 的某邻域内的闭子集上，可见
\[
  (\cO^k_{\omega,\tau})_z=(R^k_{\omega,\tau,u_\tau})_{\mathfrak m_\tau}
  =(R^k_{\omega,\tau,\sigma_{u_\tau}})_{\mathfrak m_\tau},
\]
其中 $\mathfrak m_\tau$ 是同样由幺半群理想 $P\setminus\{0\}$ 生成的极大理想。这里最后一个等式成立，是因为 $R^k_{\omega,\tau,u_\tau}$ 是 $R^k_{\omega,\tau,\sigma_{u_\tau}}$ 在某个不属于 $\mathfrak m_\tau$ 的元素处的局部化。为避免记号过于繁复，引入简写
\[
  R_\tau:=(R^k_{\omega,\tau,\sigma_{u_\tau}})_{\mathfrak m_\tau}.
\]
现在需要证明存在同构
\[
  \varprojlim R_\tau\xrightarrow{\ \cong\ }\bigl(\Bbbk[P]/(t^{k+1})\bigr)_{\mathfrak m_k}.
\]
左端的逆向系统由映射
\[
  \varphi_{\tau_1,\tau_2}:R_{\tau_2}\to R_{\tau_1}
\]
给出，它们是 $F_k(\epsilon_{\tau_1,\tau_2})$ 的局部化。这是换层映射与若干换腔映射的复合。我们可以通过在 $\sigma$ 中到 $v$ 的平行移动，把各种 $\sigma$ 对应的幺半群 $P_{\varphi,\omega,\sigma}$（其中 $\omega=\{v\}$）与 $P_{\varphi,v}$ 等同。这给出典范的换层映射
\[
  \psi_{\tau_1,\tau_2}:R_{\tau_2}\to R_{\tau_1},
\]
于是 $\varphi_{\tau_1,\tau_2}$ 是 $\psi_{\tau_1,\tau_2}$ 与一列换腔映射
\[
  \theta_{u_{n-1},u_n}\circ\cdots\circ\theta_{u_1,u_2}:R_{\tau_1}\to R_{\tau_1}
\]
的复合，后者用 $\omega$ 附近的点 $y$ 定义（故 $B$ 的奇点在此不起作用）。注意环的集合 $\{R_\tau\}$ 连同映射 $\varphi_{\tau_1,\tau_2}$ 按通常方式构成逆向系统，

% [原文 p.289]
而同一组环连同映射 $\psi_{\tau_1,\tau_2}$ 也构成逆向系统。后者的逆向极限由引理 6.9 给出，即 $(\Bbbk[P]/(t^{k+1}))_{\mathfrak m_k}$。因此只需给出这两个逆向系统之间的同构，即自同构 $\psi_\tau:R_\tau\to R_\tau$，使得只要 $\omega\subseteq\tau_1\subseteq\tau_2$，就有
\begin{equation}
  \begin{array}{ccc}
    R_{\tau_2} & \xrightarrow{\ \psi_{\tau_2}\ } & R_{\tau_2}\\[2pt]
    {\scriptstyle\varphi_{\tau_1,\tau_2}}\big\downarrow & & \big\downarrow{\scriptstyle\psi_{\tau_1,\tau_2}}\\[2pt]
    R_{\tau_1} & \xrightarrow{\ \psi_{\tau_1}\ } & R_{\tau_1}
  \end{array} \tag{6.5}
\end{equation}
交换。

为此，首先注意若 $u,u'$ 是相邻的腔，则可写
\[
  \theta_{u,u'}(z^m)=z^mf_{u,u'}^{\langle n,\bar m\rangle}
\]
对某个 $n$ 成立，且 $f_{u,u'}$ 可看作由结构 $\cS$ 完全确定的 $\Bbbk[P]$ 中某元素的像。此外 $f_{u,u'}\notin\mathfrak m_\tau$，故对每个 $\tau$，$f_{u,u'}$ 可看作 $R_\tau^\times$ 的元素。于是可以如下定义
\[
  \theta^\tau_{u,u'}:R_\tau\to R_\tau,
\]
即令
\[
  \theta^\tau_{u,u'}(z^m)=z^mf_{u,u'}^{\langle n,\bar m\rangle}.
\]
注意 $\psi_{\tau_1,\tau_2}\circ\theta^{\tau_2}_{u,u'}=\theta^{\tau_1}_{u,u'}\circ\psi_{\tau_1,\tau_2}$。

现在对任意 $\tau\supseteq\omega$，令 $u_\tau=u_1,u_2,\ldots,u_n=u_\omega$ 为从 $u_\tau$ 过渡到 $u_\omega$ 所得的腔的序列。若 $\omega\not\subseteq\partial B$，我们约定以顺时针方式进行，使得该链良定义。然后定义
\[
  \psi_\tau=\theta^\tau_{u_{n-1},u_n}\circ\cdots\circ\theta^\tau_{u_1,u_2}.
\]
容易验证，如此选取后，(6.5) 交换。

% [Tutor Note: (6.5) 交换图中 $\psi_\tau$ 的构造要求“顺时针”约定（当 $\omega\not\subseteq\partial B$ 时），否则链不唯一、$\psi_\tau$ 会差一个换腔映射；这与 §6.2.3 中单值变换的计算是同一回事。核对锚点：$f_{u,u'}\notin\mathfrak m_\tau$ 保证 $\theta^\tau_{u,u'}$ 是自同构（$f_{u,u'}\in R_\tau^\times$）。]

该引理现在表明 (6.4) 在 $z$ 处正合，从而如上所述 (6.3) 正合。于是 $U^k_\omega$ 是 $U^0_\omega$ 的平坦形变。

我们只需证明它在点 $z$ 处 étale 同构于 $\Bbbk[P]/(t^{k+1})$。这由上面的引理可得。那里给出的同构诱导出映射
\[
  \psi:\Bbbk[P]/(t^{k+1})\to(R^k_\omega)_{\mathfrak m_z},
\]
其中 $\mathfrak m_z$ 是点 $z$ 在 $R^k_\omega$ 中的极大理想。由于 $P$ 是有限生成的，可以找到 $f\in R^k_\omega\setminus\mathfrak m_z$，使对每个 $m\in P$ 有 $\psi(z^m)\in(R^k_\omega)_f$。于是得到映射
\[
  \psi:\Bbbk[P]/(t^{k+1})\to(R^k_\omega)_f,
\]
它在 $\mathfrak m_k$ 与 $\mathfrak m_z$ 处局部化后成为同构，从而给出概形映射 $V\to\operatorname{Spec}\Bbbk[P]/(t^{k+1})$，它在 $z$ 处 étale。由 étale 性质的开性（例如见 [83] 的命题 3.8），得到所需结果。

为构造 $X_0=\check X_0(B,\cP)$ 的形变 $X_k$，我们沿开子集粘合。确实，若 $\omega_1\subseteq\omega_2$，则典范地有 $U^k_{\omega_2}\subseteq U^k_{\omega_1}$。这一典范包含如下得到。在定义 $U^k_{\omega_2}$ 时，需要对每个 $\tau\supseteq\omega_2$ 选取参考腔 $u^2_\tau$。类似地，对每个 $\tau\supseteq\omega_1$ 选取参考腔 $u^1_\tau$。于是若 $\tau\supseteq\omega_2$，就得到良定义的映射
\[
  R^k_{\omega_2,\tau,u^2_\tau}\to R^k_{\omega_1,\tau,u^1_\tau},
\]

% [原文 p.290]
它是换腔同态与换层同态的复合。这诱导出逆向极限上的映射，从而给出映射 $U^k_{\omega_2}\to U^k_{\omega_1}$，容易看出它是开集的包含。细节留给读者。

构造了这些典范包含之后，我们可以沿这些公共开集把诸概形 $U^k_\omega$ 粘合起来，得到 $X_k$。这些直接的细节同样留给读者。于是我们证明了

\begin{theorem}
若 $\cS$ 是到 $k$ 阶相容的结构，则存在在 $O_k$ 上平坦的概形 $\check X_k(B,\cP)$，它是 $\check X_0(B,\cP)$ 的形变。
\end{theorem}

\begin{example}
继续例 6.23，我们希望描述原点的顶点 $\omega$ 所对应的概形 $U^k_\omega$。读者可以验证 $R^k_\omega$ 的如下描述：使用变量 $x_1,\ldots,x_5$，它们应看作 $z^{(2,1,0)}$、$z^{(3,1,1)}$、$z^{(2,0,1)}$、$z^{(0,-1,0)}$、$z^{(0,0,-1)}$ 在映射 $R^k_\omega\to R^0_\omega$ 下的提升。这里最后两个坐标位于 $\Lambda$ 在 $\omega$ 处的茎中，而第一个坐标给出 $\varphi_\omega$ 的值。则
\[
  R^k_\omega\cong\Bbbk[x_1,\ldots,x_5,t]/(t^{k+1},f_1,\ldots,f_5),
\]
其中
\[
  \begin{aligned}
    f_1&=x_5x_2-t(x_1+t^2)\\
    f_2&=x_1x_3-t(x_2+t^3)\\
    f_3&=x_2x_4-t(x_3+t^2)\\
    f_4&=x_3x_5-t^2(x_4+1)\\
    f_5&=x_4x_1-t^2(x_5+1).
  \end{aligned}
\]
为把它紧化，不必在对应于 $\cP$ 的另外五个顶点的邻域中寻找图表，只需把上面的图表射影化：增加一个变量，并在 $\mathbb{A}^1\times\PP^5$ 中得到由方程
\[
  \begin{aligned}
    x_5x_2-t(x_1x_0+t^2x_0^2)&=0\\
    x_1x_3-t(x_2x_0+t^3x_0^2)&=0\\
    x_2x_4-t(x_3x_0+t^2x_0^2)&=0\\
    x_3x_5-t^2(x_4x_0+x_0^2)&=0\\
    x_4x_1-t^2(x_5x_0+x_0^2)&=0
  \end{aligned}
\]
给出的 del Pezzo 曲面的退化族。
\end{example}

\subsection{实现相容性：热带顶点群}

剩下要做的一切就是构造到 $k$ 阶相容的结构 $\cS$，然后令 $k\to\infty$。为此，我们需要理解射线相撞时会发生什么。

\subsubsection{Kontsevich 与 Soibelman 的论证}

我们首先描述一个相当一般的框架与一般的结果。这将提供产生相容结构的基本想法。这里的定义只是 §§5.4.2 与 5.4.3 中所见定义的轻微推广，因此应当看起来熟悉。

% [原文 p.291]
固定格 $M\cong\Z^2$，$N=\operatorname{Hom}_\Z(M,\Z)$ 如常，并设给定带映射
\[
  r:P\to M
\]
的环面幺半群 $P$。本章的典型例子是 $P=P_{\varphi,x}$（$x\in B_0$），而 $r:P_{\varphi,x}\to M$ 是由 $m\mapsto\bar m$ 给出的投影 $P_{\varphi,x}\to\Lambda_x$。然后我们定义 $\Bbbk[P]$ 的（相对）\emph{对数导子}模为
\[
  \Theta(\Bbbk[P]):=\operatorname{Hom}_\Z(M,\Bbbk[P])=\Bbbk[P]\otimes_\Z N.
\]
元素 $f\otimes n$ 记作 $f\partial_n$，它作为 $\Bbbk[P]$ 上的导子作用为
\[
  f\partial_n(z^m)=f\langle n,r(m)\rangle z^m.
\]
现在用 $\mathfrak m$ 同时表示理想 $P\setminus P^\times$（其中 $P^\times$ 记 $P$ 的单位群）与它在 $\Bbbk[P]$ 中生成的单项式理想。例如若 $P=M\oplus\mathbb{N}$，$r$ 为到 $M$ 的投影，则 $P^\times=M\oplus0$。设 $I$ 是 $\Bbbk[P]$ 中根为 $\mathfrak m$ 的单项式理想；我们称这样的理想为 \emph{$\mathfrak m$-准素}理想。于是给定 $\xi\in\mathfrak m\Theta(\Bbbk[P])$，通过导子 $\xi$ 的指数化得到元素
\[
  \exp(\xi)\in\operatorname{Aut}(\Bbbk[P]/I),
\]
即
\[
  \exp(\xi)(a)=a+\sum_{i=1}^{\infty}\frac{\xi^i(a)}{i!}.
\]
这是有限和，因为存在整数 $n$ 使 $\mathfrak m^n\subseteq I$。

例如可以验证，若 $f\in\mathfrak m$ 可写成 $f=\sum_{m\in S}c_mz^m$，其中 $\{r(m)\mid m\in S\}$ 落在 $M$ 的一个秩一子格中，且 $n\in N$ 零化该子格，则
\[
  \exp(f\partial_n)(z^m)=\exp(f)^{\langle n,r(m)\rangle}z^m.
\]
然而一般地，计算向量场的指数是困难的。

对数导子模上有一个自然的 Lie 括号：
\[
  \begin{aligned}
    [z^m\partial_n,z^{m'}\partial_{n'}]
    &:=z^{m+m'}\bigl(\langle n,r(m')\rangle\partial_{n'}
    -\langle n',r(m)\rangle\partial_n\bigr)\\
    &=z^{m+m'}\partial_{\langle n,r(m')\rangle n'-\langle n',r(m)\rangle n}.
  \end{aligned}
\]
现在令
\[
  \mathfrak v:=\bigoplus_{\substack{m\in\mathfrak m\\ r(m)\ne0}}
  z^m\Bbbk\otimes r(m)^\perp\subseteq\Theta(\Bbbk[P]).
\]
这是由满足 $m\in\mathfrak m$、$r(m)\ne0$、$\langle n,r(m)\rangle=0$ 的向量场 $z^m\partial_n$ 的线性组合构成的 $\Bbbk$-向量空间。通过与 §5.4.2 中完全相同的论证，$\mathfrak v$ 在 Lie 括号下封闭。对每个 $\mathfrak m$-准素理想 $I$，于是得到 $\operatorname{Aut}(\Bbbk[P]/I)$ 的子群
\[
  \mathbb{V}_I:=\{\exp(\xi)\mid\xi\in\mathfrak v\}.
\]
正与 §5.4.2 中一样，我们可以把该群的元素表示成 Hamilton 辛同构。然而，这一点在本章中对我们将不重要。

% [Tutor Note: 本节的 $\mathfrak v$、$\mathbb{V}_I$ 就是 §5.4.2 中“热带顶点群”的雏形：$\mathfrak v$ 在 Lie 括号下封闭的论证与 §5.4.2 完全相同（按 $r(m)$ 的秩一性质分情形验证）。核对锚点：条件 $\langle n,r(m)\rangle=0$ 保证 $\partial_n$ 沿 $m$ 的方向不产生线性项，从而 $\exp(\xi)$ 对每个 $\mathfrak m$-准素理想 $I$ 都良定义在 $\operatorname{Aut}(\Bbbk[P]/I)$ 中。]

有时我们希望在一个完全的情形下处理幂级数。令
\[
  \widehat{\Bbbk[P]}:=\varprojlim\Bbbk[P]/\mathfrak m^k
\]
为 $\Bbbk[P]$ 在理想 $\mathfrak m$ 处的完备化。

% [Tutor Note: 本块末句按 STYLE_GUIDE §10「跨页句子开始者写完」补完：原文 p.291 末的 “Let $\widehat{\Bbbk[P]}:=\varprojlim\Bbbk[P]/\mathfrak m^k$” 在 p.292 首行接 “be the completion of $\Bbbk[P]$ at the ideal $\mathfrak m$.”。下一块 c06d 应从 “We can then also define $\widehat{\mathbb{V}}:=\varprojlim\mathbb{V}_{\mathfrak m^k}$, a pro-nilpotent subgroup of $\operatorname{Aut}(\widehat{\Bbbk[P]})$.” 起译。]
